\documentclass[11pt]{article}
\usepackage{amsmath, amssymb}

\title{Approximation of Mutual Information in Random Walk Voltage Models}
\author{}
\date{}

\begin{document}

\maketitle

\section*{Problem Setting}

Consider a graph with $n$ nodes. We designate $k$ source nodes $S_1, \ldots, S_k$ and a single sink node $T$. For each source $S_i$, define the voltage function $V_i(u)$ on nodes $u \in \{1, \ldots, n\}$ via the Dirichlet problem:

\[
\begin{cases}
V_i(S_i) = 1, \\
V_i(T) = 0, \\
\Delta V_i(u) = 0 \text{ for } u \notin \{S_i, T\}.
\end{cases}
\]

Each $V_i(u)$ represents the probability that a random walk starting at node $u$ reaches source $S_i$ before the sink $T$.

\bigskip

Suppose $N$ independent particles are released from node $u$. For each source $S_i$, let $m_i(u)$ be the number of particles that hit $S_i$ before $T$. Then:

\[
m_i(u) \sim \operatorname{Binomial}(N, V_i(u)).
\]

Define the $k$-dimensional random vector at node $u$ as:

\[
A(u) = \big(m_1(u), m_2(u), \ldots, m_k(u)\big).
\]

We introduce a random variable $B$ that is uniformly distributed over the $n$ nodes:

\[
B \sim \text{Uniform} \{1, 2, \ldots, n\}.
\]

We are interested in the mutual information $I(A; B)$, where $A = A(B)$.

\section*{Conditional Distribution of $A$}

Conditioned on $B = u$, the components of $A$ are independent:

\[
A \mid B = u \sim \bigotimes_{i=1}^k \operatorname{Binomial}(N, V_i(u)).
\]

For large $N$, each binomial distribution can be approximated by a Gaussian:

\[
\operatorname{Binomial}(N, p) \approx \mathcal{N}\big(Np, Np(1-p)\big).
\]

Therefore,

\[
A \mid B = u \approx \mathcal{N}\big(\mu(u), \Sigma(u)\big),
\]

where

\[
\mu_i(u) = N V_i(u),
\]
\[
\Sigma(u) = \operatorname{diag}\big(N V_1(u)(1 - V_1(u)), \ldots, N V_k(u)(1 - V_k(u))\big).
\]

\section*{Mutual Information Formula}

The mutual information is:

\[
I(A; B) = H(A) - H(A|B).
\]

\subsection*{Conditional Entropy $H(A|B)$}

Given $B = u$, $A$ is approximately Gaussian with covariance $\Sigma(u)$. The differential entropy of a multivariate Gaussian is:

\[
H(A|B = u) = \frac{k}{2}\log(2\pi e) + \frac{1}{2} \log \det \Sigma(u).
\]

Averaging over uniform $B$:

\[
H(A|B) = \frac{1}{n} \sum_{u=1}^n H(A|B = u).
\]

\subsection*{Marginal Entropy $H(A)$}

The marginal distribution of $A$ is a mixture of Gaussians:

\[
P(A) = \frac{1}{n} \sum_{u=1}^n \mathcal{N}\big(\mu(u), \Sigma(u)\big).
\]

Its entropy $H(A)$ has no closed-form expression. We approximate it via Monte Carlo sampling:

\begin{itemize}
    \item Sample $B$ uniformly from $\{1, \ldots, n\}$.
    \item Sample $A \mid B$ from $\mathcal{N}\big(\mu(B), \Sigma(B)\big)$.
    \item Estimate $p(A)$ using the mixture density.
    \item Approximate $H(A) \approx - \frac{1}{M} \sum_{j=1}^M \log p(A^{(j)})$.
\end{itemize}

\section*{Monte Carlo Approximation Procedure}

\begin{enumerate}
    \item For each node $u$, compute:
    \[
    \mu(u) = N \cdot V(u), \quad \Sigma(u) = \operatorname{diag}\big(N V_i(u)(1 - V_i(u))\big).
    \]
    
    \item Compute:
    \[
    H(A|B) = \frac{1}{n} \sum_{u=1}^n \frac{k}{2}\log(2\pi e) + \frac{1}{2} \sum_{i=1}^k \log\big(N V_i(u)(1 - V_i(u))\big).
    \]
    
    \item Draw $M$ samples:
    \[
    B^{(j)} \sim \text{Uniform}\{1, \ldots, n\},
    \]
    \[
    A^{(j)} \sim \mathcal{N}\big(\mu(B^{(j)}), \Sigma(B^{(j)})\big).
    \]
    
    \item For each sample $A^{(j)}$, compute:
    \[
    p(A^{(j)}) \approx \frac{1}{n} \sum_{u=1}^n \mathcal{N}\big(A^{(j)}; \mu(u), \Sigma(u)\big).
    \]
    
    \item Approximate:
    \[
    H(A) \approx -\frac{1}{M} \sum_{j=1}^M \log p(A^{(j)}).
    \]
    
    \item Finally:
    \[
    I(A; B) \approx H(A) - H(A|B).
    \]
\end{enumerate}

\section*{Interpretation}

\begin{itemize}
    \item For large $N$, the binomial distributions are sharply peaked, so the means of the Gaussians separate well.
    \item If the voltage vectors $V(u)$ are well-separated across nodes, $I(A; B)$ approaches $\log n$.
    \item If the voltage vectors are similar across nodes, $I(A; B)$ is small.
\end{itemize}

This approximation provides a practical way to quantify how well observing the particle count vector $A$ can identify the underlying node $B$ in large-$N$ regimes.

\end{document}

%%% Local Variables:
%%% mode: latex
%%% TeX-master: t
%%% End:
